\section{The parity theorem}\label{sec:parity}

Fix $n=2r+1$ and a reduced word $w=i_1\cdots i_N$ of $w_0\in S_n$. Wire $x$ starts at height $x$
(heights $1,\dots,n$ from top to bottom). Letter $i_t$ is the \emph{crossing} $t$: the wires at heights
$i_t$ and $i_t+1$ exchange heights. Along each wire, its $n-1$ crossings are numbered $1,\dots,n-1$ in
time order. These are its \emph{positions}.

\begin{definition}
A crossing is an \emph{$A$-crossing} if it has an even position along both of its wires.
Equivalently, both wires have made an odd number of crossings before it. Let $\#A(w)$ be the number
of $A$-crossings.
\end{definition}

In $G$ every edge joins an odd position to the next even position on one wire. So an $A$-crossing is
a vertex both of whose $G$-neighbours occur earlier in time.

\begin{lemma}\label{lem:percycle}
Every cycle of $G$ contains an odd number of $A$-crossings. Consequently the number of
components of $G$ is congruent to $\#A(w)$ modulo $2$.
\end{lemma}

\begin{proof}
Process the crossings in time order. After $t$ steps, call a wire \emph{open} if it has made an odd
number of crossings. Its last crossing is then waiting for its $G$-partner along that wire.

Let $G_t$ be the subgraph of $G$ formed by the processed crossings and the edges with both ends
processed. Every component of $G_t$ is a cycle of $G$ or a path. Each path has two \emph{dangling
ends}, and each dangling end belongs to an open wire; distinct open wires carry distinct dangling
ends. A path consisting of a single vertex has both dangling ends at that vertex.

Pairing the two dangling ends of each path gives a perfect matching $\mu_t$ on the set $S_t$ of heights
currently occupied by open wires. Consider crossing $t+1$, between wire $u$ at height $i$ and wire $v$
at height $i+1$, and call it $\tau$. Exactly one of the following happens.
\begin{itemize}
\item[(B)] \emph{Both wires closed.} $\tau$ has an odd position on both wires. It starts a new path
  whose dangling ends now sit at heights $i,i+1$, so the pair $(i,i+1)$ is added to $\mu$.
\item[(C)] \emph{Exactly one wire open.} Say $u$ is open; the case of $v$ is symmetric. Then $\tau$
  extends the path that dangled on $u$. Its new dangling end is on $v$, which now occupies height $i$,
  the height $u$ had. So $S$ and $\mu$ are unchanged.
\item[(A)] \emph{Both wires open.} $\tau$ is an $A$-crossing and joins the dangling ends at heights $i$
  and $i+1$. If $\mu(i)=i+1$ this \emph{closes} a cycle of $G$. Otherwise it \emph{joins} two paths:
  remove $(i,\mu(i))$ and $(i+1,\mu(i+1))$ and add $(\mu(i),\mu(i+1))$.
\end{itemize}
Wires not involved in $\tau$ do not move.

By induction on $t$ we prove two invariants.
\begin{itemize}
\item[(I1)] $\mu_t$ is non-crossing for the order of heights. Case (B) adds a pair of adjacent heights,
  which cannot interleave with any pair. In case (A) the pairs $(i,\mu(i))$ and $(i+1,\mu(i+1))$ do not
  cross and $i,i+1$ are adjacent. So they lie either side by side,
  $\mu(i)<i<i+1<\mu(i+1)$, or nested, $i<i+1<\mu(i+1)<\mu(i)$ or its mirror image. A pair crossing the
  merged pair would cross one of the two old pairs.
\item[(I2)] For every pair $(a,b)\in\mu_t$ with $a<b$, the difference $b-a$ is odd. This is true for
  $(i,i+1)$. Under a side-by-side join, $b-a=(i-a)+1+(b-i-1)$ is a sum of two odd numbers plus $1$.
  Under a nested join, the new difference is (odd)$-$(odd)$-1$.
\end{itemize}

Fix a cycle $C$ of $G$ and let $\Lambda_C(t)=\sum (b-a-1)/2$, the sum over the pairs of $\mu_t$ whose
path lies in $C$. By (I2) this is an integer. We compare $\Lambda_C$ before and after each step that
involves $C$:
\begin{itemize}
\item a (B)-step adds $0$;
\item a (C)-step changes nothing;
\item closing removes $(i,i+1)$, so it changes $\Lambda_C$ by $0$;
\item a side-by-side join changes it by $\bigl((b-a-1)-(i-a-1)-(b-i-2)\bigr)/2=1$;
\item a nested join $(i,a)\supset(i+1,b)$ changes it by $-(b-i-1)$, which is odd by (I2);
\item the mirror nested join $(b,i+1)\supset(a,i)$ changes it by $a-i$, again odd.
\end{itemize}
Steps at crossings not on $C$ do not touch the pairs of $C$.

The earliest vertex of $C$ has both neighbours later in time, so it is a (B)-vertex and
$\Lambda_C=0$ right after it. Just before $C$ closes, the only pair of $C$ is $(i,i+1)$, so again
$\Lambda_C=0$. Each closing step closes exactly one cycle, so $C$ has exactly one closing step.
Therefore $C$ has an even number of joins, and its $A$-crossings (joins plus one closing) are odd in
number. Summing over the cycles gives the congruence.
\end{proof}

\begin{lemma}\label{lem:A}
For every reduced word $w$ of $w_0\in S_{2r+1}$ we have $\#A(w)\equiv r(r+1)/2\pmod 2$.
\end{lemma}

\begin{proof}
A wire changes height by $\pm1$ at each of its crossings. So the number of crossings it has made so
far is congruent modulo $2$ to its current height minus its label.

Consider a crossing at heights $(i,i+1)$ of wires $u$ and $v$. They have not crossed before, so their
initial order persists and $u<v$. The crossing is an $A$-crossing iff $i-u$ and $i+1-v$ are both odd.
If $u+v$ is even this is impossible. If $u+v$ is odd it holds iff $i-u$ is odd.

\emph{Commutation moves.} Exchanging adjacent letters $i_t,i_{t+1}$ with $|i_t-i_{t+1}|\ge2$ changes
neither the wires nor the heights of any crossing.

\emph{Braid moves.} Replacing $i,i+1,i$ by $i+1,i,i+1$ involves three wires $x<y<z$. Each of the three
crossings $xy,xz,yz$ keeps its wires, and its upper height changes by exactly $1$. So the move flips the
$A$-status of exactly those of the three pairs whose labels have odd sum. Among three integers the
number of such pairs is $0$ or $2$, so $\#A$ changes by an even number.

\emph{Connectivity.} By the theorem of Matsumoto and Tits~\cite{Matsumoto1964}, any two reduced words
of $w_0$ are connected by commutation and braid moves. Hence $\#A \bmod 2$ is the same for all reduced
words of $w_0$.

\emph{Computation.} Take the bubble-sort word $(s_1s_2\cdots s_{n-1})(s_1\cdots s_{n-2})\cdots(s_1)$.
Here every wire meets the other wires in increasing order of labels. For $u<v$ the crossing therefore
has position $v-1$ on $u$ and position $u$ on $v$. It is an $A$-crossing iff $u$ is even and $v$ is
odd. Hence
\[
  \#A=\sum_{j=1}^{r}\#\{v\ \text{odd}: 2j<v\le 2r+1\}=\sum_{j=1}^{r}(r-j+1)=\frac{r(r+1)}2. \qedhere
\]
\end{proof}

\begin{proof}[Proof of Theorem~\ref{thm:parity}]
Combine Lemmas~\ref{lem:percycle} and~\ref{lem:A}, and note that $(n^2-1)/8=r(r+1)/2$.
\end{proof}

\begin{proof}[Proof of Corollary~\ref{cor:main}]
For $n=2r+1$, the number $r(r+1)/2$ is even iff $r\equiv0,3\pmod 4$, i.e.\ iff $n\equiv\pm1\pmod 8$.
In that case $G$ has an even number of components, so it is not a single cycle. By
Proposition~\ref{prop:forced}(1), no simple $\binom n2$-gon lies on $n$ lines, and by Proposition~\ref{prop:forced}(2) every
such polygon has at most $\binom n2-3$ sides.
\end{proof}

\begin{remark}
For $n\equiv\pm3\pmod 8$, Theorem~\ref{thm:parity} only says that $G$ has an odd number of components. It does
not guarantee a single cycle; for example, both $1$ and $3$ components occur for $n=5$.
\end{remark}
